% !PRINTED 313
\weberpage{313}
%<<CONT
confine ourselves to the first cases, and so we begin with the quadratic equations
\weq{1}{x^2 + a\,x + b = 0.}
If we regard $a$ and $b$ as rectangular coordinates in a plane, then every point of this plane is the image of one and only one quadratic equation, and conversely every quadratic equation with real coefficients has a point of the plane as its image.

The equation
\weq{2}{a^2 - 4\,b = 0}
represents a parabola, on which lie the images of the equations with equal roots. The inner points, that is, those lying in the hollow of the parabola, represent the equations with imaginary roots, the outer points the equations with real roots.

We now form, according to \S.\,89, the \gesperrt{Sturm} chain
\weq{3}{x^2 + a\,x + b,\quad 2\,x + a,\quad a^2 - 4\,b}
and examine the changes of sign. If $x$ is constant and $a$, $b$ variable, then
\[
x^2 + a\,x + b = 0
\]
is the equation of the tangent to the parabola at the point
\[
a = -2\,x,\quad b = x^2,
\]
and indeed $x^2 + a\,x + b$ is positive in the part of the plane that also contains the parabola.

By this tangent and by the parabola the plane is divided into four regions, in which, as is indicated in Fig. 16 (on the following page), 0, 1, 2 changes of sign occur in the sequence (3).

If we now set $x = \alpha$ and $x = \beta$, we obtain two such tangents and a division of the plane into four regions, in which lie the representative points of those equations that have no root, one root, or two roots between $\alpha$ and $\beta$. In Fig. 17 these four regions are marked by the corresponding numerals.

This is the difference between those points from which no tangent, one tangent, or two tangents can be drawn to the parabola having their point of contact between $\alpha$ and $\beta$. Fig. 17 thus gives the exact and complete distinction of the various cases.
